\subsection*{BR-008: the permanent P-20 witness has exact border rank eighteen}
\textbf{Status: proved.  Author complete exact rank computations, root whole-proof
and all-new-source review, and independent root coefficient/stabilizer/cover
checks have passed.  The root wrapper run explicitly did not recompute ranks.
This statement concerns one ten-plane only.}

Let $S_{20}=\{0,1,2,3,4,6,7,8,9,10\}$, and let $T_{20}$ be the
literal integral permanent residual tensor for this coordinate ten-plane.
Over an algebraically closed field of characteristic zero,
\[
 \underline R(T_{20})=18.
\]
This does not determine the global permanent minimum, which remains
$17\le\mu_{10}(U^{\mathrm{per}}_{4,2})\le18$.

\paragraph{A connected solvable stabilizer.}
Write $T_{20}=H\boxtimes G$, where
\[
 H=\alpha_1\otimes(r_0\otimes r_1+r_1\otimes r_0)
  +\alpha_2\otimes(r_0\otimes r_2+r_2\otimes r_0),
\]
and $G$ has the five symmetric off-diagonal slices indexed by
$01,02,03,12,13$ on the four column coordinates.  The product first
factor is ten-dimensional, and the last two factors have coordinates
$r_i\otimes c_j$ for $0\le i\le2$, $0\le j\le3$.
The full row/column diagonal torus stabilizes this tensor after the
contragredient action on slice labels.  We use four additional
one-parameter unipotent subgroups:
\begin{enumerate}
\item For $i=1,2$, act on the row coordinate in the second factor by
$r_i\mapsto r_i+s r_0$, and in the third factor by
$r_i\mapsto r_i-s r_0$.  Keep the first factor fixed.  The two new
$r_0\otimes r_0$ terms in the $i$th star slice cancel, and all other
slices are unchanged.
\item Act in both last factors by $r_2\mapsto r_2+s r_1$, with the
contragredient action on the two row slice labels.  The row star
space is preserved.
\item Act in both last factors by $c_3\mapsto c_3+s c_2$, with the
contragredient action on the five column slice labels.  The slices
$03,13$ map respectively to themselves plus $s$ times $02,12$;
all remaining slices are unchanged.
\end{enumerate}
These subgroups and the diagonal torus generate a connected solvable
algebraic group: the two opposite-translation subgroups form a
normal additive plane under the common upper-triangular row-leaf
group, while the column-leaf upper-triangular group commutes with
the row actions.  Equivalently the non-torus part is generated by
three row upper-triangular root operators and one independent
column root, with their bracket still in the two-dimensional
translation span.  Thus the generated group is a quotient of a
product of connected solvable groups.  This argument uses no
assertion that the determinant's full Borel group preserves a
permanent tensor.

The checker verifies directly the zero torus weight of each of the
forty nonzero tensor coefficients and the vanishing of each of the
four infinitesimal full three-factor actions.  In characteristic zero
these infinitesimal identities integrate to the stated unipotent
identities.  On $A\otimes B$ the two translation roots act only on
$B$; the other two roots have the displayed dual action on $A$.
These are exactly the four operators in
\texttt{probe\_per10\_p20.py}.

\paragraph{Complete flag and pencil cover.}
If $\underline R(T_{20})\le17$, the proper Grassmannian incidence
argument of BR-004 supplies a seventeen-plane
$E\supset T_{20}(C^*)$ whose tautological tensor has border rank at
most seventeen.  The closed incidence is invariant under the
connected solvable group just described.  Its fixed-point theorem
therefore supplies an invariant such $E$.  Since
$\dim T_{20}(C^*)=12$, its quotient has dimension five.  Lie's theorem
gives a complete invariant flag in that quotient.  The quotient
ambient module has dimension $120-12=108$, in seventy torus weight
blocks.  Within each weight, all possible next lines are the full
joint kernel of the four roots modulo the preceding flag, exactly
as in BR-004.  A joint kernel of dimension two is retained as its
entire projective line; no generic point is substituted for it.

Here is the complete rigid part of the enumeration.  The final column
counts projective families cumulatively; the generated and remaining
counts concern the rigid branches only.
\[
\begin{array}{c|r|r|r}
\text{quotient level}&\text{generated}&\text{remaining}&\text{families}\\\hline
1&11&11&0\\
2&76&71&5\\
3&402&222&58\\
4&1368&73&297\\
5&708&0&392
\end{array}
\]
All three hundred ninety-two exported families are projective lines.
Their first levels are respectively $2,3,4,5$ with counts
$5,53,239,95$.  There is no higher-dimensional exported parameter
space and no full rigid survivor.

The exact univariate minor procedure of BR-005 is applied to every
one of these lines.  It computes a polynomial minor, factors its
possible zero set, evaluates every irreducible factor in its exact
quotient field, and checks infinity.  It uses the full higher Koszul
maps and, when needed, the fixed seven-dimensional projected map,
whose pure-tensor cost is twenty and whose threshold is $340$.
Three hundred fifty initial lines are excluded by the full Koszul
pencils without continuation.  Ten further lines are excluded by
the projected test without continuation.  Thirty-two lines retain
finite points and are continued by the same complete joint-kernel
flag procedure.  One of those thirty-two finite reductions also
uses the projected pencil.  The only points requiring flag
continuation are $t=0$, $t=1$, and infinity.  Each subsequent
projective line is again subjected to the full exact pencil test;
all finite successors are inserted, and any unhandled family would
be retained explicitly.  Every continuation terminates with no
full flag and no open family.  Thus the final certificate contains
exactly $392$ excluded initial families, zero unhandled families,
and zero full survivors.  All possible invariant $E$ have been
excluded, contradicting border rank at most seventeen.

\paragraph{Upper bound and reproduction.}
The upper bound uses the elementary integral first- and second-difference
curves in predecessor P-20, \texttt{sections/per10\_upper18.tex},
especially its displayed $H(t)$ and $G(t)$ and their Kronecker product.
They have three and six pure terms respectively.  The independent
standard-library routine \texttt{check\_per10\_p20.py} collects all
coefficients of their eighteen-term product.  It finds no negative
power, and its constant coefficients agree with all forty literal
permanent tensor entries.  It also checks the full stabilizer
identities above and validates the lower cover's index completeness.
Its optional \texttt{--rebuild} mode regenerates the flag cover and
every polynomial-rank computation using the already reviewed
standard-library/SymPy engines; without that option it explicitly
records that ranks were not recomputed.  The initial complete flag
run took $68.7802$ seconds and the complete refinement took
$204.3852$ seconds.  The exact machine-readable traces are
\texttt{per10\_p20\_gate.json}, \texttt{per10\_p20\_refined.json},
and \texttt{per10\_p20\_summary.json}.

\paragraph{A rejected global-screening shortcut.}
A chosen seven-slice restriction with border rank at least fifteen
only implies that the ten-slice tensor has border rank at least
fifteen.  It does not by itself add three further ranks.  The
three-step border-substitution inequality uses the minimum over
all seven-dimensional slice subspaces.  After torus reduction this
minimum runs over every coordinate seven-subset inside the given
ten-plane.  Therefore a sound hereditary test must establish the
strong lower bound on all those subsets, or prove an additional
condition forcing the endpoint into a specified strong locus.
The stronger one-subset implication was suggested during this run,
immediately rejected, and used in no theorem or elimination count.
